\begin{textsample}{POS Dim 4 – 2012\_c6 – Score 1.00 – 2012\_c6/enade\_2012\_tec\_gest\_com\_8fb0b169d0.txt}  \label{ex:f4_pos_126}
CAPÍTULO 3 ANÁLISE TÉCNICA DA PROVA Este capítulo tem por objetivo apresentar o desempenho dos estudantes concluintes de Tecnologia em Gestão Comercial no ENADE/2012.
Para isso, foram calculadas as estatísticas básicas da prova em seu todo, bem como as estatísticas dos componentes relacionadas à Formação Geral, ao de Conhecimento Específico da Área e das questões discursivas isoladamente.
Nas tabelas, são apresentados o tamanho da população inscrita e de presentes, e as seguintes estatísticas das notas⁹ : média do desempenho na prova, erro padrão da média, desvio padrão, nota mínima, mediana e nota máxima.
Tais estatísticas contemplam o total de estudantes concluintes da Área de Tecnologia em Gestão Comercial presentes à prova do ENADE 2012, tendo -se em vista as seguintes agregações : ( a ) as Grandes Regiões e o país como um todo ; ( b ) a Categoria Administrativa ; e ( c ) a Organização Acadêmica.
Em relação aos gráficos de distribuição de notas, o intervalo considerado foi de 10 unidades, aberto à esquerda e fechado à direita, com exceção do primeiro intervalo, [ 0 ; 10 ], fechado em ambos os extremos.
Para os gráficos de distribuição das notas das questões discursivas, foram consideradas mais duas categorias : questão em branco e nota zero. 3.1 ESTATÍSTICAS BÁSICAS DA PROVA 3.1.1 Estatísticas Básicas Gerais A Tabela 3.1 apresenta as estatísticas básicas da prova por grande Região.
A população total de inscritos foi de 8.553.
Destes, 6.409 estiveram presentes, sendo 25,1 \% o índice de não comparecimento.
A Região de maior abstenção foi a Sul ( 29,2 \% ) e a de menor abstenção foi a Centro-Oeste ( 17,2 \% ).
A média das notas da prova como um todo ( nas seções seguintes serão analisados os componentes de Formação Geral e de Conhecimento Específico ) foi 42,1, sendo que os alunos da região Norte obtiveram a média mais baixa ( 39,4 ) e os da região Sul obtiveram a média mais alta ( 43,4 ).
As demais médias foram : 42,7 na região Centro-Oeste, 41,6 na região Sudeste e 40,1 na região Nordeste.
O desvio padrão para o Brasil como um todo foi 14,1, sendo o maior desvio padrão encontrado na região Nordeste ( 15,7 ) e o menor na região Centro-Oeste ( 12,8 ), indicando uma menor dispersão das notas desta última região. ⁹ Essas estatísticas e outras estão definidas no Capítulo 1. 14,1, sendo o maior desvio padrão encontrado na região Nordeste ( 15,7 ) e o menor na região Centro-Oeste ( 12,8 ), indicando uma menor dispersão das notas desta última região.
A região que obteve a maior nota máxima foi a Sudeste ( 84,9 ), ao passo que a região que atingiu a menor nota máxima foi a Centro-Oeste ( 73,3 ).
A mediana do Brasil como um todo foi 42,3, sendo a maior mediana obtida na região Sudeste ( 43,5 ) e a menor obtida na Norte ( 40,0 ).
A nota mínima foi zero na maioria das regiões com exceção das regiões Norte com 1,9 e Centro-Oeste com 8,9.
O comportamento das notas dos estudantes de todo o Brasil pode ser observado no Gráfico 3.1 que apresenta um histograma com a distribuição das mesmas.
Essa é uma distribuição unimodal com moda no intervalo ( 40 ; 50 ].
Apesar do coeficiente de assimetria da distribuição das notas ser negativo ( -0,14 ) este é pequeno e podemos considerar que a distribuição é aproximadamente simétrica.
As distribuições por Grande Região também apresentam assimetria negativa na maioria, concentração pouco maior do lado direito do histograma e mais espalhada do lado esquerdo.
As únicas exceções são as regiões Norte e Sudeste, com coeficiente de assimetria positivo ( 0,08 e 0,02, respectivamente ), com o comportamento oposto : concentração à esquerda e cauda um pouco mais longa à direita.
Os Gráficos 3.2, 3.3 e 3.4 apresentam informações referentes à média da nota final dos Participantes, desagregando os resultados de acordo com, respectivamente, as Grandes Regiões do país, a Categoria Administrativa e a Organização Acadêmica.
Os gráficos apresentam o valor da média das notas como uma barra e os extremos do intervalo de confiança de 95 \% como linhas verticais unidas por uma linha horizontal na forma da letra H maiúscula.
Considerando -se o gráfico de notas segundo Grande Região ( Gráfico 3.2 ), observa -se que existe diferença estatisticamente significativa ao nível de 95 \% entre a maior média, obtida na região Sul ( 43,4 ) e a menor, obtida na região Norte ( 39,4 ).
Levando -se em conta os estudantes por Categorias Administrativas ( Gráfico 3.3 ), observa -se que existe diferença estatisticamente significante entre as médias das notas das IES Públicas e Privadas.
Pode ser observado ainda que a média dos alunos de IES Públicas ( 47,4 ) é maior que a dos alunos de IES Privadas ( 42,0 ).
A diferença entre IES Públicas e Privadas ( 5,4 ) é maior do que a diferença máxima entre regiões ( 4,0 ).
Tendo como foco o Gráfico 3.4, que apresenta as notas médias das provas segundo Organização Acadêmica, constata -se que existe diferença estatisticamente significativa ao nível de 95 \% nas médias das notas dos estudantes provenientes dos diferentes tipos de Organização Acadêmica.
A média dos estudantes provenientes de Universidades ( 40,4 ) foi menor do que a dos estudantes provenientes de Centros Universitários ( 42,2 ) e Faculdades ( 43,3 ). 3.1.2 Estatísticas Básicas no Componente de Formação Geral A Tabela 3.2 apresenta as estatísticas básicas em relação ao componente da prova que avalia a Formação Geral dos estudantes concluintes.
Os alunos de todo Brasil obtiveram desempenho médio de 41,0.
Quanto à variabilidade, o desvio padrão das notas dos estudantes do Brasil como um todo foi 15,6.
A maior média foi obtida na região Sul ( 42,4 ), e a menor, na região Norte ( 37,6 ).
As demais médias foram : 41,5 na região Centro-Oeste, 40,6 na região Sudeste e 38,2 na região Nordeste.
Já o maior desvio padrão foi obtido na região Nordeste ( 16,5 ) e o menor na região Norte ( 14,5 ).
Os demais desvios padrões foram : 15,5 na região Sul e 15,3 nas regiões Sudeste e Centro-Oeste.
A maior nota no componente de Formação Geral da prova do ENADE foi obtida por pelo menos um aluno da região Sudeste ( 88,0 ) enquanto que a menor nota máxima foi obtida na região Norte ( 64,0 ).
Nas outras regiões a nota máxima foi a mesma, 86,0.
A mediana do Brasil como um todo foi 41,5, sendo a menor mediana encontrada na região Norte ( 38,0 ) e a maior encontrada na região Sul ( 43,0 ).
A nota mínima nesta parte foi zero em quase todas as regiões, exceto na região Centro-Oeste ( 7,5 ).
O Gráfico 3.5 propicia a avaliação do desempenho dos estudantes no componente de Formação Geral a partir do histograma da distribuição das notas correspondentes.
A distribuição é unimodal, com moda em ( 40 ; 50 ], assim como na prova como um todo.
Nota -se, ainda, que no gráfico 3.5 as notas apresentam uma maior dispersão do que no Gráfico 3.1 ( distribuição das notas da prova ), confirmado pela comparação dos desvios padrões : 14,1 para a nota da prova como um todo e 15,6 para o componente de Formação Geral.
Para o componente de Formação Geral, o coeficiente de assimetria da distribuição das notas dos estudantes, como na prova como um todo, também é negativo ( -0,15 ).
Por ser um coeficiente pequeno, a distribuição apresenta uma leve concentração à direita e cauda maior à esquerda.
Em quase todas as Grandes Regiões os histogramas também possuem assimetria negativa ( entre -0,24 na região Nordeste e -0,05 na região Sudeste ), exceto a região Centro-Oeste, onde a assimetria é positiva ( 0,04 ).
Nos Gráficos 3.6, 3.7 e 3.8 são apresentadas as informações referentes ao desempenho dos Concluintes no componente de Formação Geral, em diferentes agregações : Grande Região do país, Categoria Administrativa e Organização Acadêmica.
Observa -se no Gráfico 3.6 que há diferença estatisticamente significativa entre as médias das notas no Componente de Formação Geral, segundo Grande Região do país, se comparada a região Sul, que obteve a maior média ( 42,4 ), com a região Norte, que obteve a menor média ( 37,6 ).
Vemos que o intervalo de confiança mais largo é o da região Norte ; já os intervalos mais estreitos são observados nas regiões Sudeste e Sul.
Este fato está relacionado, também, com o tamanho da população envolvida, menor na região Norte do que no Sudeste e no Sul.
No Gráfico 3.7, que representa as notas médias no Componente de Formação Geral segundo Categoria Administrativa do país, observa -se que existe diferença estatisticamente significativa entre as médias.
Assim como ocorreu para as notas da prova como um todo, para o componente de Formação Geral, os concluintes de Tecnologia em Gestão Comercial das IES Públicas obtiveram uma média maior ( 44,8 ) do que os das IES Privadas ( 40,9 ).
Considerando -se o tipo de Organização Acadêmica, nota -se, no Gráfico 3.8, uma diferença estatisticamente significativa entre as médias das notas dos alunos provenientes de Universidades em relação às médias dos provenientes de Centros Universitários e Faculdades, enquanto não houve diferença estatisticamente significativa entre Centros Universitários e Faculdades.
Nos Centros Universitários ( 42,0 ), a média é mais elevada do que em Universidades ( 39,5 ) e Faculdades ( 41,0 ). 3.1.3 Estatísticas Básicas do Componente de Conhecimento Específico A Tabela 3.3 apresenta as estatísticas básicas referentes ao componente de Conhecimento Específico da área de Tecnologia em Gestão Comercial.
A média do desempenho dos alunos do Brasil como um todo foi 42,5.
A maior média foi obtida na região Sul ( 43,7 ), e a menor, na região Norte ( 40,0 ).
As demais médias foram : 43,0 na região Centro-Oeste, 41,9 na região Sudeste e 40,8 na região Nordeste.
Quanto à variabilidade das notas, o desvio padrão do Brasil como um todo foi 15,8, sendo o maior desvio padrão observado na região Nordeste ( 17,5 ) e o menor na região Centro-Oeste ( 14,6 ).
Os demais desvios foram : 16,6 da região Norte, 15,7 da região Sudeste e 15,5 da região Sul.
A mediana das notas dos estudantes de todo o Brasil foi 42,5.
A maior mediana ocorreu na região Centro-Oeste ( 44,0 ) e a menor na região Norte ( 38,2 ).
As demais medianas foram 43,7 na região Sul, 41,6 na região Sudeste e 41,1 na região Nordeste.
A nota máxima do Brasil como um todo foi 90,2, sendo obtida por pelo menos um aluno da região Sudeste.
As demais notas máximas foram : 89,8 na região Sul, 88,3 na região Nordeste, 88,3 na região Nordeste, 84,6 na região Norte e 80,4 na região Centro-Oeste.
A nota mínima foi zero em quase todas as regiões, com exceção da região Centro-Oeste, cuja nota mínima foi 6,8.
Assim como os Gráficos 3.1 e 3.5, o Gráfico 3.9, apresentado a seguir, proporciona uma avaliação do desempenho de concluintes em relação ao componente de Conhecimento Específico com um histograma da distribuição das notas correspondentes.
Assim como as outras duas distribuições apresentadas, esta é aproximadamente simétrica.
A distribuição é unimodal, e o grupo modal também é o ( 40 ; 50 ].
O coeficiente de assimetria da distribuição das notas do componente de Conhecimento Específico é negativo e bem próximo de zero ( -0,04 ).
Já nas regiões Norte e Sudeste, o coeficiente de assimetria é positivo, iguais a 0,18 e 0,06, respectivamente, evidenciando que a distribuição das notas dos estudantes dessas regiões, no componente de Conhecimento Específico, tem cauda mais pesada à direita.
Os Gráficos 3.10, 3.11 e 3.12 apresentam, respectivamente, uma comparação dos resultados em relação à Grande Região do país, à Categoria Administrativa e à Organização Acadêmica, agora levando em conta o desempenho de alunos no componente de Conhecimento Específico da prova.
Pelo Gráfico 3.10, observa -se dois grupos, dentro dos quais não existe diferença estatisticamente significativa entre as médias das notas no componente de Conhecimento Específico, um formado pelas regiões Centro-Oeste, Sul e Sudeste, e o outro formado pelas regiões Nordeste e Norte.
Sendo que a região Sul ( 43,7 ) foi a que obteve média mais elevada, enquanto a região Norte ( 40,0 ) foi a região que obteve média mais baixa.
Quanto à Categoria Administrativa ( Gráfico 3.11 ), observa -se um comportamento semelhante àquele da parte de Formação Geral e à prova como um todo, ou seja, existe diferença estatisticamente significativa entre as médias das IES Públicas e Privadas, sendo que a maior média foi obtida por alunos de IES Públicas de ensino ( 48,2 ).
Quanto ao Gráfico 3.12, observa -se, mais uma vez, que existe diferença estatisticamente significativa ao nível de 95 \% entre as notas no componente de Conhecimento Específico dos diferentes tipos de Organização Acadêmica.
Sendo que a média dos Concluintes das Faculdades ( 44,0 ) foi maior do que de Centros Universitários ( 42,2 ) e de Universidades ( 40,7 ). 3.2 ANÁLISE DAS QUESTÕES OBJETIVAS 3.2.1 Componente de Formação Geral A Tabela 3.4 apresenta as estatísticas básicas relativas às oito questões objetivas do componente da prova que abrange a Formação Geral dos estudantes.
A média do Brasil foi 44,4.
A menor média foi encontrada na região Norte ( 39,3 ) e a maior na região Sul ( 46,0 ).
As demais médias foram 44,4 na região Centro-Oeste, 44,0 na região Sudeste e 41,6 na região Nordeste.
O desvio padrão do Brasil foi 19,8, sendo o maior desvio padrão encontrado na região Nordeste ( 20,2 ) e o menor na região Norte ( 17,8 ).
Os demais desvios foram : 19,7 nas regiões Sul e Centro-Oeste, e 19,5 na região Sudeste.
A mediana foi 50,0 para quase todas as regiões do Brasil, exceto para as regiões Norte e Nordeste ( 37,5 ).
As notas máximas, nas questões objetivas do componente de Formação Geral, foram 100,0 em quase todas as regiões, com exceção da região Norte, cuja máxima foi 75,0.
Já as notas mínimas ( 0,0 ) foram iguais para todas as regiões, sem exceção.
A Tabela 3.5 apresenta os índices de facilidade e discriminação ( ponto bisserial ) das questões objetivas do componente de Formação Geral.
Quanto ao índice de facilidade, foram usadas as seguintes cores para diferenciar o nível de dificuldade da questão : • Azul para as questões classificadas com índice muito fácil ( índice \&gt ; =0,86 ), verde para as questões classificadas com índice fácil ( 0,61 a 0,85 ), amarelo para as questões classificadas com médio ( 0,41 a 0,60 ), vermelho para as questões classificadas com difícil ( 0,16 a 0,40 ) e roxo para as questões classificadas com muito difícil ( \&lt ; =0,15 ).
Já quanto ao índice de discriminação, foram usadas as seguintes cores para qualificar a questão : • As questões classificadas com índice fraco receberam a cor roxa ( índice \&lt ; =0,19 ), as classificadas com médio receberam a cor vermelho ( 0,20 a 0,29 ), as classificadas com bom receberam a cor verde ( 0,30 a 0,39 ) e as classificadas com muito bom ( \&gt ; =0,40 ) receberam a cor azul.
As questões objetivas do componente de Formação Geral, segundo o índice de facilidade, foram assim avaliadas : seis questões foram consideradas de dificuldade média, situando -se no intervalo entre 0,41 e 0,60 do índice de facilidade, ou seja, houve entre 41,0 \% e 60,0 \% de acertos ; as duas outras questões foram classificadas como difíceis, situando -se no intervalo entre 0,16 e 0,40.
Como já comentado, para análise das questões objetivas relativas à Formação Geral segundo o poder de discriminação, utilizou -se, o índice de discriminação ponto bisserial.
Nesta análise as questões foram assim avaliadas : cinco das oito questões apresentaram 3.2 ANÁLISE DAS QUESTÕES OBJETIVAS 3.2.1 Componente de Formação Geral A Tabela 3.4 apresenta as estatísticas básicas relativas às oito questões objetivas do componente da prova que abrange a Formação Geral dos estudantes.
A média do Brasil foi 44,4.
A menor média foi encontrada na região Norte ( 39,3 ) e a maior na região Sul ( 46,0 ).
As demais médias foram 44,4 na região Centro-Oeste, 44,0 na região Sudeste e 41,6 na região Nordeste.
O desvio padrão do Brasil foi 19,8, sendo o maior desvio padrão encontrado na região Nordeste ( 20,2 ) e o menor na região Norte ( 17,8 ).
Os demais desvios foram : 19,7 nas regiões Sul e Centro-Oeste, e 19,5 na região Sudeste.
A mediana foi 50,0 para quase todas as regiões do Brasil, exceto para as regiões Norte e Nordeste ( 37,5 ).
As notas máximas, nas questões objetivas do componente de Formação Geral, foram 100,0 em quase todas as regiões, com exceção da região Norte, cuja máxima foi 75,0.
Já as notas mínimas ( 0,0 ) foram iguais para todas as regiões, sem exceção.
A Tabela 3.5 apresenta os índices de facilidade e discriminação ( ponto bisserial ) das questões objetivas do componente de Formação Geral.
Quanto ao índice de facilidade, foram usadas as seguintes cores para diferenciar o nível de dificuldade da questão : • Azul para as questões classificadas com índice muito fácil ( índice \&gt ; =0,86 ), verde para as questões classificadas com índice fácil ( 0,61 a 0,85 ), amarelo para as questões classificadas com médio ( 0,41 a 0,60 ), vermelho para as questões classificadas com difícil ( 0,16 a 0,40 ) e roxo para as questões classificadas com muito difícil ( \&lt ; =0,15 ).
Já quanto ao índice de discriminação, foram usadas as seguintes cores para qualificar a questão : • As questões classificadas com índice fraco receberam a cor roxa ( índice \&lt ; =0,19 ), as classificadas com médio receberam a cor vermelho ( 0,20 a 0,29 ), as classificadas com bom receberam a cor verde ( 0,30 a 0,39 ) e as classificadas com muito bom ( \&gt ; =0,40 ) receberam a cor azul.
As questões objetivas do componente de Formação Geral, segundo o índice de facilidade, foram assim avaliadas : seis questões foram consideradas de dificuldade média, situando -se no intervalo entre 0,41 e 0,60 do índice de facilidade, ou seja, houve entre 41,0 \% e 60,0 \% de acertos ; as duas outras questões foram classificadas como difíceis, situando -se no intervalo entre 0,16 e 0,40.
Como já comentado, para análise das questões objetivas relativas à Formação Geral segundo o poder de discriminação, utilizou -se, o índice de discriminação ponto bisserial.
Nesta análise as questões foram assim avaliadas : cinco das oito questões apresentaram O Gráfico 3.13, para exemplificar, analisa o comportamento da questão de número 7 de Formação Geral.
Trata -se da quarta questão mais fácil e a que obteve o maior índice de discriminação dessa parte da prova.
Neste gráfico, cada uma das cinco curvas representa o percentual de respostas em determinada alternativa da questão, em função da nota dos estudantes nesta parte da prova ( Formação Geral/Múltipla Escolha ), antes de possíveis eliminações pelo critério do ponto bisserial.
A curva em verde corresponde à alternativa B, a correta para esta questão.
Assim, observa -se que entre os estudantes com notas mais baixas, nessa parte do exame, a situação mais frequente foi a escolha de uma das alternativas incorretas : a alternativa C ( em preto ) ou D ( em roxo ).
Na medida em que a nota aumenta, indicando desempenho melhor índices acima de 0,40 e, assim, foram classificadas com índice muito bom para esse grupo de alunos ; as outras três questões tiveram bom índice de discriminação, entre 0,30 e 0,39.
O índice de facilidade variou de 0,32 a 0,55, e o de discriminação, de 0,30 a 0,45.
As questões com índices de discriminação muito bom, de números 1, 2, 4, 6 e 7, figuraram entre as mais fáceis desse conjunto : todas classificadas na categoria médio.
Em particular, as questões 2 e 6 apresentaram igual poder discriminatório, com índice 0,42, sendo que a questão 6 foi mais fácil, com uma proporção de 0,51 de acertos, enquanto a questão 2 teve proporção de 0,42.
A questão de número 5 apresentou índice de facilidade 0,32, ou seja, um quantitativo de 32,0 \% dos estudantes conseguiu resolvê -la, dentro do universo de participantes.
Além disso, seu índice de discriminação, 0,30, foi o menor entre as questões.
A questão 8 obteve índice de discriminação bom, 0,35, e seu índice de facilidade foi difícil.
Já a questão de número 3, considerada média quanto à facilidade, obteve índice de discriminação 0,39, considerado bom.
O Gráfico 3.13, para exemplificar, analisa o comportamento da questão de número 7 de Formação Geral.
Trata -se da quarta questão mais fácil e a que obteve o maior índice de discriminação dessa parte da prova.
Neste gráfico, cada uma das cinco curvas representa o percentual de respostas em determinada alternativa da questão, em função da nota dos estudantes nesta parte da prova ( Formação Geral/Múltipla Escolha ), antes de possíveis eliminações pelo critério do ponto bisserial.
A curva em verde corresponde à alternativa B, a correta para esta questão.
Assim, observa -se que entre os estudantes com notas mais baixas, nessa parte do exame, a situação mais frequente foi a escolha de uma das alternativas incorretas : a alternativa C ( em preto ) ou D ( em roxo ).
Na medida em que a nota aumenta, indicando desempenho melhor nesta parte da prova, aumenta concomitantemente a proporção de estudantes que selecionaram a alternativa correta B, atingindo 100 \% para os alunos com 8 acertos.
Essa análise permite verificar como a questão discriminou os grupos de desempenho, justificando o alto índice obtido na questão.
Os gráficos relativos às demais questões de Formação Geral constam do Anexo I. 3.2.2 Componente de Conhecimento Específico A Tabela 3.6 apresenta as estatísticas básicas em relação às questões objetivas do componente de Conhecimento Específico da prova, por Grande Região.
A média do Brasil deste componente foi de 42,4.
A menor média foi observada na região Norte ( 40,3 ) e a maior na região Sul ( 43,7 ).
O desvio padrão de todo o Brasil foi 16,4, sendo o menor desvio padrão encontrado na região Centro-Oeste ( 15,5 ) e o maior na região Nordeste ( 18,2 ).
A mediana de todo o Brasil foi 44,4, sendo a mesma encontrada nas regiões Sul e Centro-Oeste.
As demais medianas foram : 41,7 na região Norte e 38,9 nas regiões Nordeste e Sudeste.
A nota máxima da prova ( 94,4 ) foi obtida, nas questões objetivas do componente de Conhecimento Específico, por pelo menos um aluno da região Sudeste.
A nota máxima nas demais regiões foi de 88,9 nas regiões Nordeste e Sul e 83,3 nas regiões Norte e Centro-Oeste.
A nota mínima foi zero em quase todas as regiões do Brasil, com exceção da região Centro-Oeste ( 5,6 ).
A Tabela 3.7 apresenta os índices de facilidade e discriminação ( ponto bisserial ) das questões objetivas do componente de Conhecimento Específico.
Para facilitar a diferenciação das questões usaram -se as mesmas cores da Tabela 3.5 para as diferentes classificações dos índices de facilidade e de discriminação.
A partir dos índices obtidos das questões objetivas da parte da prova relativa ao componente de Conhecimento Específico, pode -se concluir que a maioria das questões foi considerada pelo menos difícil : das 27 questões, onze foram classificadas como difíceis e cinco como muito difíceis.
Não houve questão classificada como muito fácil, ao passo que uma foi tida como fácil, na faixa de 0,61 a 0,85 do índice de facilidade, e outras dez consideradas médias, entre 0,41 e 0,60.
Já quanto aos índices de discriminação das questões objetivas do componente de Conhecimento Específico da prova, tem -se como resultado a seguinte classificação : oito das 27 questões foram consideradas como boas, enquanto três delas tiveram índice de discriminação muito bom.
Assim, para 11 das 27 questões, os índices de discriminação foram bons ou muito bons.
Dentre as demais, sete delas foram classificadas como médias e outras nove como fracas, sendo dezesseis, por conseguinte, a quantidade de questões nos dois patamares mais baixos de discriminação.
Constata -se, assim, que a prova – no que se refere ao componente de Conhecimento Específico – possuía capacidade baixa de discriminar entre aqueles que dominam ou não o conteúdo.
Dentre as questões que alcançaram os maiores índices de discriminação, as de números 28, 31 e 33 classificadas com índice muito bom, com índice acima ou igual a 0,40, duas delas ( questões 28 e 31 ) foram classificadas na categoria média e a outra como fácil, quanto ao índice de facilidade.
A questão de número 11 foi a mais difícil dentre as 27 questões específicas, com baixo índice de facilidade, apenas 7,0 \% de acertos.
Essa questão apresentou poder discriminatório igualmente baixo, 0,09, o que comprova ter sido esta a mais difícil para os estudantes.
Destaca -se, também, a questão 20, com índice de facilidade 0,08, o que, em termos percentuais, corresponde a 8,0 \% de estudantes que responderam acertadamente, obtendo, ainda, 0,04 de índice de discriminação.
Tais questões foram, portanto, duas das mais difíceis da prova.
Pelo critério do ponto bisserial, elas foram eliminadas do cômputo da nota final.
Além destas duas, as demais questões com índice fraco de discriminação, questões 10, 17, 18, 21, 22, 26 e 27 também não foram computadas.
A título de exemplo das análises do comportamento das questões objetivas, o Gráfico 3.14 analisa a questão 28 do componente de Conhecimento Específico.
Esta foi uma das questões mais fáceis da prova, apresentando índice de facilidade 0,49, ou seja, 49,0 \% dos estudantes assinalaram acertadamente a opção A, correspondente ao gabarito.
Seu índice de discriminação foi igual a 0,46, classificado como muito bom.
Neste gráfico, cada uma das cinco curvas representa o percentual de respostas em determinada alternativa da questão 28, em função da nota dos estudantes nesta parte da prova, antes de possíveis eliminações de questões pelo critério do ponto bisserial.
A alternativa correta A, representada no gráfico pela curva em azul, foi escolhida em maiores proporções pelos alunos com desempenho melhor nesta parte da prova.
Já as alternativas incorretas, também denominadas distratores, foram selecionadas principalmente por aqueles com notas mais baixas.
Aqueles com nota zero, na sua quase totalidade deixaram esta questão em branco ou marcaram mais de uma alternativa, comportamento considerado inválido.
A proporção de alunos que selecionou a resposta correta A aumenta gradativamente, chegando a atingir 100 \% para alunos com 18 acertos ou mais, enquanto a proporção dos que escolheram alternativas incorretas decai, a partir da primeira nota não nula, como função da nota nesta parte da prova.
Os gráficos relativos às demais questões do componente de Conhecimento Específico constam do Anexo I. 3.3 ANÁLISE DAS QUESTÕES DISCURSIVAS 3.3.1 Componente de Formação Geral As análises dos resultados de desempenho dos estudantes de Tecnologia em Gestão Comercial nas duas questões discursivas relativas à Formação Geral encontram -se na Tabela 3.8 e no Gráfico 3.15.
Na tabela 3.8 observa -se que as notas médias foram mais baixas nesse conjunto de questões do que no das objetivas.
Os estudantes de todo o Brasil obtiveram, em Formação Geral, média 44,4 nas questões objetivas e 35,9 nas questões discursivas.
A mediana também confirma o pior desempenho dos alunos de todo o Brasil nas questões discursivas do componente de Formação Geral.
Enquanto essa estatística foi de 37,5 para as questões discursivas, para as questões objetivas essa estatística foi de 50,0.
Pode -se, também, notar um leve aumento do desvio padrão de 19,8, nas questões objetivas do componente de Formação Geral dos alunos de todo o Brasil, para 20,1 nas questões discursivas do mesmo componente.
Como anteriormente dito, a mediana de todo o Brasil, neste componente, foi 37,5, sendo a mesma mediana encontrada nas regiões Norte, Nordeste e Sudeste, enquanto a mediana das regiões Sul e Centro-Oeste foi 40,0.
A maior nota máxima ( 85,0 ) foi obtida, nas questões discursivas do componente de Formação Geral, por pelo menos um aluno da região Sul, enquanto a menor nota máxima ( 77,5 ) foi obtida nas regiões Norte e Nordeste.
A nota mínima foi zero em todas as regiões do Brasil, sem exceção.
O Gráfico 3.15 representa a distribuição das notas nas questões discursivas no componente de Formação Geral.
A moda desta distribuição ocorre no intervalo, ( 40 ; 50 ].
Observa -se também grande quantidade de notas zero, pela alta frequência de alunos que deixaram este tipo de questão em branco, e frequência próxima da modal no intervalo ( 30 ; 40 ].
A distribuição possui assimetria à esquerda, com coeficiente de assimetria -0,35.
Em todas as regiões o coeficiente de assimetria também é negativo, o maior em módulo, -0,43, na região Sul, e o menor na região Nordeste ( -0,24 ).
Na sequência, os resultados verificados para cada uma das questões discursivas de Formação Geral serão apresentados, estabelecendo -se relações com os conteúdos abordados em cada uma delas.
Os comentários da Banca de docentes corretores a respeito do observado na correção das respostas dos estudantes, suas impressões e conclusões serão apresentados junto à análise de cada questão.
Cumpre esclarecer que, tendo em vista que as questões discursivas de Formação Geral são padronizadas, ou seja, constam de todas as provas, os comentários da Banca são os mesmos para todas as carreiras acadêmicas, sendo direcionados a todos os estudantes que participaram do ENADE/2012.
A seguir, serão analisados os desempenhos da Área de Tecnologia em Gestão Comercial nas duas questões discursivas de Formação Geral do ENADE/2012, comparando os resultados obtidos com comentários para cada questão. 3.3.1.1 Análise da Questão Discursiva 1 do Componente de Formação Geral Os dados de Tecnologia em Gestão Comercial, obtidos a partir das respostas à questão 1, encontram -se na Tabela 3.9 e no Gráfico 3.16.
Nessa questão – de pior desempenho dentre as duas de Formação Geral – os alunos de todo Brasil tiveram média, 29,6.
A maior média para a questão 1 foi obtida na região Centro-Oeste ( 31,0 ), e a menor, na região Nordeste ( 27,1 ).
Quanto à variabilidade das notas, o desvio padrão de todo o Brasil foi 23,5.
O menor desvio padrão foi obtido na região Sul ( 23,0 ) e o maior desvio padrão foi obtido na região Centro-Oeste ( 25,2 ).
As medianas e as notas mínimas da questão discursiva 1 foram as mesmas para todas as regiões do Brasil, respectivamente, 30,0 e 0,0.
Além disso, a maior nota máxima ( 95,0 ) foi obtida nas regiões Sudeste e Sul do país, enquanto a menor nota máxima ( 80,0 ) foi obtida na região Norte.
O Gráfico 3.16 mostra a distribuição das notas na questão discursiva 1 do componente de Formação Geral.
Observa -se que a maior frequência, correspondendo à moda, é na categoria dos alunos que deixaram a questão em branco.
Considerando apenas os estudantes que tentaram resolver a questão, a moda ocorre no intervalo ( 20 ; 30 ].
A distribuição possui assimetria à direita, com coeficiente de assimetria positivo ( 0,26 ).
Em todas as regiões do Brasil este coeficiente se mantém positivo ( entre 0,21 na região Sul e 0,32 nas regiões Norte e Nordeste ). 3.3.1.2 Comentários sobre a correção das respostas à Questão Discursiva 1 A questão aborda um tema bastante polêmico o qual vem sendo amplamente discutido, a conciliação entre o desenvolvimento econômico e a preservação dos recursos naturais.
Busca -se observar a relação entre o estímulo governamental à compra de veículos novos e os impactos ambientais decorrentes desta política pública.
Sua compreensão implica no conhecimento de uma temática atual e na reflexão critica a respeito das decisões governamentais em sua correlação com o cenário econômico internacional e os problemas e soluções de ordem ambiental.
Embora o comando da questão solicitasse explicitamente que a resposta fosse fornecida em um texto único ( “ redija um texto dissertativo sobre... ” ), foi significativamente alta a proporção de respondentes que atendeu a cada item ( a, b e c ) do enunciado de maneira isolada, como partes estanques da mesma resposta.
No que se segue encontra -se uma avaliação da correção de cada um dos três itens. a - “ conceito de desenvolvimento sustentável ” A variedade de respostas a esse item confirmou a importância do tema para avaliar a formação geral dos estudantes e cursos, já que ( i ) é atual ; ( ii ) foi abordado nos meios de comunicação ( especialmente no ano de 2012, em que o Brasil sediou a “ Rio+20 ” ) ; ( iii ) é relevante não apenas no Brasil como em todos os principais fóruns internacionais ; e ( iv ) é comum a diversas áreas de conhecimento.
Outros dois aspectos positivos merecem destaque : ( i ) a grande incidência de respostas que apontaram a relação entre gerações como aspecto essencial da ideia de sustentabilidade socioambiental – no mais, as respostas apresentaram pelo menos a tentativa de indicar uma definição de desenvolvimento sustentável ou a abordagem positiva dessa mesma noção, evidenciado compreensão do enunciado ; e ( ii ) o exercício crítico que a questão propunha, no sentido de exigir dos respondentes a compreensão das contradições de políticas públicas e a apresentação de propostas de soluções para conciliar o incentivo do crescimento econômico e a promoção da sustentabilidade socioambiental, a partir de uma perspectiva transformadora da realidade.
Apesar dos aspectos positivos apontados no parágrafo anterior, a maioria dos estudantes não alcançou o máximo da pontuação.
De forma geral, trataram o desenvolvimento sustentável como sinônimo de não agressão ambiental.
Alguns confundiram desenvolvimento sustentável com desenvolvimento econômico, e outros faziam referência apenas à preocupação com a conservação dos recursos mencionando a preocupação com gerações futuras.
Deve -se, ainda, destacar o significativo número de alunos que apresentou uma definição aparentemente copiada do enunciado de outra questão da prova, como por exemplo, a questão objetiva 31 da prova de Administração. b - “ conflito entre o estímulo à compra de veículos automotores e a promoção da sustentabilidade ” Esse foi o item identificado como aquele no qual as respostas dos estudantes avaliados mais se aproximaram da pontuação máxima prevista no padrão de resposta.
Ou seja, a ampla maioria dos respondentes foi bem-sucedida em apontar a relação entre o aumento da quantidade de veículos automotores em circulação e os danos socioambientais que isso acarreta no cenário urbano, devido ao aumento da emissão de gases poluentes.
No item b, as respostas corretas se concentraram em críticas ao incentivo oferecido pelo governo, referências ao problema dos congestionamentos e ao aumento da poluição.
Contudo, a maioria das respostas não explicitava a correlação entre o aumento da circulação de veículos automotores e o agravamento dos problemas ambientais.
Um número expressivo de alunos apresentou uma resposta em que era dado destaque à política de redução do IPI como um mecanismo para renovação da frota de automóveis ou para a aquisição do transporte próprio e da consequente redução da dependência ao transporte coletivo.
Também foram notados alguns tipos de padrões recorrentes nas tentativas de resposta a este item.
Dois desses padrões se destacaram : ( i ) as declarações genéricas de condenação do “ consumismo ” de uma “ nova classe média ” ou da busca pelo “ status ” que a posse de um automóvel particular proporcionaria, como se fosse possível atribuir a essa ordem de coisas a causa principal do aumento dos congestionamentos nas grandes cidades ; e ( ii ) a repetição da avaliação, presente no enunciado da questão, que aponta de forma superficial a contradição entre o estímulo à demanda por automóveis particulares e a promoção da sustentabilidade. c - “ ações de fomento ao transporte urbano sustentável no Brasil ” Este foi o item que ofereceu maior facilidade de aplicação do padrão de resposta, mas foram poucos os casos de pontuação integral.
Dentre as soluções mais destacadas pelos alunos apareceram : ( i ) o incentivo ao uso de energias alternativas ; e ( ii ) o estímulo ao transporte coletivo ( metrô, trem e ônibus ), além do uso de bicicletas.
Importante destacar o fato de que muitos apresentavam respostas com referências à necessidade de incentivo ao transporte urbano, nas quais este termo era utilizado como sinônimo de transporte coletivo.
A avaliação da questão 1 trouxe à tona problemas estruturais da educação básica, que se refletem no ensino superior.
Foi muito grande o número de provas em que o aluno demonstrou não compreender o enunciado da questão, apresentando respostas que fugiam ao tema proposto, mesmo quando bem redigidas.
Uma resposta recorrente fazia referência às consequências da ampla oferta de crédito à população e aos problemas relacionados à inadimplência, situação que, embora, verdadeira, não estava relacionada à questão proposta.
Houve, ainda, os casos em que as respostas apresentadas não só se afastaram do tema como não reuniam qualquer tipo de argumentação coerente.
E, por fim, destaca -se a existência de respostas incoerentes e desconexas.
Outro aspecto geral diz respeito ao grande número de provas em branco ou em que o aluno apresentou a resposta a outro item no espaço da questão 1.
Considerando -se o objetivo geral deste exame, avaliar a qualidade dos cursos de graduação oferecidos no país e contribuir para a melhoria da qualidade do ensino superior, destaca -se a imensa responsabilidade do processo que envolve a elaboração deste diagnóstico.
Nesse sentido, é urgente alertar para a ocorrência de um número significativo de respostas que se afastam de um padrão mediano de qualidade esperado de um aluno concluinte em um curso de graduação, especialmente em se considerando o fato de que a questão proposta faz referência a tema intensamente explorado pela mídia. 3.3.1.3 Análise da Questão Discursiva 2 do Componente de Formação Geral A Tabela 3.10 mostra que o desempenho dos estudantes na questão 2 ( média 42,3 ) foi tipicamente superior ao obtido na questão de número 1 ( média 29,6 ).
A região Sul foi aquela onde a média, nessa questão, foi maior ( 43,6 ), e a de menor média foi a região Nordeste ( 39,3 ).
Quanto à variabilidade das notas, o desvio padrão de todo o Brasil foi 23,6, superior ao obtido na questão de número 1 ( 23,5 ).
O maior desvio nessa questão foi obtido na região Nordeste ( 25,4 ), enquanto o menor foi obtido na região Norte ( 21,6 ).
A mediana foi a mesma nas regiões Norte e Nordeste ( 45,0 ), enquanto a mediana das demais regiões foi igual a 50,0, tal qual a do Brasil como um todo.
A maior nota máxima ( 90,0 ) foi obtida nas regiões Sudeste e Sul, enquanto a menor nota máxima ( 80,0 ) foi obtida na região Norte.
A nota mínima foi zero em todas as regiões do Brasil, sem exceção.
O Gráfico 3.17 mostra a distribuição das notas na questão discursiva 2 do componente de Formação Geral.
Neste gráfico, observa -se grande quantidade de alunos que deixaram a questão 2 em branco ( em torno de 15 \% ), e a moda no intervalo ( 40 ; 50 ].
O coeficiente de assimetria das notas nesta questão é negativo ( -0,56 ), indicando que a distribuição possui assimetria à esquerda.
O mesmo ocorre para todas as regiões ( entre -0,63 na região Sul e -0,41 na região Nordeste ). 3.3.1.4 Comentários sobre a correção das respostas à Questão Discursiva 2 O conteúdo da questão era bastante adequado às Diretrizes do INEP para o ENADE, além de atual e de abrangente.
Porém, exigia -se a articulação de muitas variáveis.
Buscando aferir, a partir da apresentação de uma charge e de um texto introdutório, as competências e habilidades relacionadas à leitura, à interpretação, à análise crítica, à capacidade argumentativa, à elaboração de sínteses e de avaliação prospectiva sobre tema da atualidade, a questão tratava da violência e suas repercussões no contexto atual.
O comando da questão solicitava, a partir da definição de violência da OMS e de duas imagens, que o aluno elaborasse um texto, de caráter dissertativo, sobre a “ violência na atualidade ”.
Embora tal solicitação fosse explícita ( “ redija um texto dissertativo sobre... ” ), como na questão 1, foi alta a proporção de alunos que respondeu a cada item do enunciado isoladamente, como partes independentes.
Verificaram -se poucas respostas em branco, como também foram poucas as respostas contendo protestos contra o ENADE, muito frequentes em anos anteriores, o que significa um avanço para a afirmação do ENADE como um exame aceito socialmente. a -"tecnologia e violência"Tanto as charges quanto o texto apresentado no enunciado induziram os alunos a focarem suas respostas nos meios de comunicação.
Além disso, algumas respostas estabeleciam a tecnologia, em si, como causa da violência tanto nas escolas ( em particular ), quanto na sociedade ( no geral ).
Muitos respondentes tiveram dificuldade para fazer a relação entre tecnologia, violência e a escola, de modo adequado, centrando -se na violência nas escolas.
Além disso, a grande maioria das respostas girava em torno da exposição da violência através da TV, internet e jogos digitais, sendo que predominaram as abordagens com ênfase em programas e filmes veiculados pela televisão.
Poucos foram os estudantes que estabeleceram relação entre o contexto socioeconômico e a violência. b -"causas e consequências da violência na escola"Embora a solicitação desse item fosse objetiva e atendida por grande parte dos alunos, inúmeras vezes, observou -se confusão entre o significado de causa e de consequência.
Em outras inúmeras vezes, as respostas apontavam apenas para uma das classes, causas ou consequências.
Foram frequentes as respostas que atribuíam o aumento da violência à ineficiência do poder público, apontando a falta de investimentos, o baixo salário dos professores, a ausência de profissionais qualificados como causas.
Foram frequentes também os que apontavam a violência doméstica, o bullying, e a transferência do imaginário dos filmes e jogos para a realidade, como causas da violência. c -"proposta de solução para o problema da violência na escola"Sem dúvida, há uma grande diversidade de abordagens possíveis para apresentação de uma proposta de solução para o problema da violência, o que o padrão de resposta buscou contemplar.
No entanto, houve grande número de respostas baseadas em opiniões pessoais dos estudantes e vinculadas ao senso comum.
As respostas giraram em torno da necessidade de mais investimentos na educação por parte do poder público e seu rebatimento no maior aparelhamento das escolas, salário dos professores etc.
Além disso, foram muitas as respostas que trataram da necessidade de maior conscientização dos pais, alunos e professores com a realização de um trabalho conjunto que envolvesse todos os segmentos da sociedade.
De um modo geral, os estudantes demonstraram dificuldades em relação ao nível de abrangência exigido, falta de conteúdo, inabilidade na elaboração de texto escrito.
Da grande quantidade de provas corrigidas, o que ficou evidente, foi a dificuldade de estabelecer as conexões requeridas na questão 2.
A violência nas escolas foi bastante associada ao reflexo do meio onde a criança e o adolescente mantém os seus vínculos afetivos.
Surgiram muitas respostas revelando críticas à intolerância no núcleo familiar, à jornada de trabalho dos pais, ao distanciamento da mãe nas atividades domésticas etc.
Em boa parte das respostas, os estudantes lançaram a causa da violência escolar exclusivamente na ausência de participação familiar na vida dos filhos, às vezes criando uma cisão com a proposta da questão, já que abandonavam totalmente a temática da tecnologia em suas reflexões.
Além da ausência dos pais no cotidiano da escola, divórcios e atos de violência contra crianças foram considerados fatores que delegariam para a escola novas funções.
Também foram mencionadas a má formação de professores e a insatisfação com o ambiente escolar.
Em inúmeros casos, ficou evidente o anseio de se fazer da escola um lugar de resgate tanto da criança e do adolescente quanto da família, compreendendo -se a escola em sentido social e assistencial.
Segundo este amplo grupo de concluintes, a escola deveria oferecer -se como centro de cultura e lazer, como espaço para o acompanhamento psicológico e serviço social às comunidades diretamente envolvidas.
Ressalte -se o constante aparecimento das sugestões de horário integral, melhores salários para professores e iniciativas para melhoria na formação de professores.
As propostas de solução para acabar com a violência na escola se detiveram em aspectos muito gerais e poucos alunos detalharam uma proposta de intervenção, apresentando ações concretas.
Dentre as propostas de solução surgiram : educação familiar, ensino religioso, maior segurança nas escolas, com utilização de aparelhos para detecção de armas, e o retorno do papel do professor como um profissional respeitado pelos poderes municipais e estaduais.
Um percentual significativo de provas apresentou resposta simplista, com ênfase em determinados pontos como : religião ( sem demonstração do conceito de laicidade ), políticas sociais ( sem demonstração do domínio de conceitos que envolvem Estado, Governo e Politicas Públicas ) etc.
Foi grande a incidência de respostas que se referiam às leis, à vigilância e maior controle dos pais e professores ; no entanto, não demonstram entender a noção de legalidade para os diferentes segmentos sociais e como os mesmos se encontram submetidos a diferentes formas de poder. noção de legalidade para os diferentes segmentos sociais e como os mesmos se encontram submetidos a diferentes formas de poder.
Alguns alunos, ainda que raros, trataram a questão com a profundidade esperada de concluintes de cursos universitários.
Por outro lado, foram vários os textos que apresentaram uma argumentação simplória e desconexa, com forte visão paternalista sobre o papel do governo e do professor, como se tais instâncias pudessem intervir na educação familiar.
Enquanto alguns conseguiram desenvolver o tema proposto com elegância e desenvoltura, encadeando o tema de forma harmônica, com excelente nível de linguagem escrita e correção gramatical, outros sequer conseguiram conectar as ideias e construir um texto razoável, utilizando, inclusive, um linguajar mais propício à linguagem falada do que à escrita.
Cabe também ressaltar a quantidade de respostas carregadas de conteúdo com juízo de valor discriminatório ou extremamente conservador, do tipo : “ mulher não deve trabalhar fora ” ; “ amar e temer a Deus é a solução ” ; obrigatoriedade do ensino religioso e de uma doutrina religiosa ; imposição de uma educação do tipo militar nas escolas ; censura nos meios de comunicação etc.
Outras soluções apontadas para resolver o problema da violência nas escolas estavam associadas à, por exemplo : redução da maioridade penal, defesa da pena de morte, milícias armadas na porta das escolas, punições como o retorno da palmatória.
O Estatuto da Criança e do Adolescente foi uma peça legal criticada, e recebeu, em diversas respostas, o título de responsável pelo sentimento de impunidade de adolescentes infratores.
Contudo, algumas dessas proposições, desde que bem contextualizadas e formuladas com boa argumentação, foram pontuadas de forma a respeitar o direito democrático de pluralidade opinativa.
É imperioso ressaltar que a falta de inteligibilidade segue representando um aspecto merecedor de providências na formação universitária, uma vez que exprime uma habilidade importante para qualquer profissional de nível superior.
Além de terem sido poucas respostas boas e claras, foi possível identificar a deficiência de argumentação, além dos inúmeros problemas gramaticais ( ortografia, concordância, pontuação, regência, etc. ).
Salienta -se, nesse sentido, o uso indiscriminado de letras maiúsculas e minúsculas numa mesma frase, sem obedecer ao critério formal da língua portuguesa. 3.3.2 Componente de Conhecimento Específico Na parte da prova relativa às questões discursivas no componente de Conhecimento Específico ( Tabela 3.11 ), observa -se que a média foi mais alta do que para as questões discursivas do componente de Formação Geral.
Enquanto no componente de Formação Geral a média para estudantes de Tecnologia em Gestão Comercial de todo o Brasil foi 35,9, na parte de Conhecimento Específico a média foi 42,7.
A maior média deste componente foi obtida pelos estudantes das regiões Sul e Centro-Oeste ( 44,1 ), e a menor, pelos da região Nordeste ( 37,7 ).
Quanto à variabilidade das notas, o desvio padrão de todo o Brasil foi 29,0.
O maior desvio padrão foi encontrado na região Nordeste ( 31,0 ), e o menor, na região Norte ( 26,5 ).
A nota máxima do país ( 95,0 ), nas questões discursivas do componente de Conhecimento Específico, foi obtida em quase todas as regiões, com exceção da região Norte, suja máxima foi 91,7.
Além disso, a nota mínima ( 0,0 ) foi obtida por alunos de todas as regiões do Brasil.
A mediana do Brasil como um todo foi 48,3.
As medianas das Grandes Regiões foram : 50,0 nas regiões Sudeste, Sul e Centro-Oeste, 40,0 na região Nordeste, e 38,4 na região Norte.
O Gráfico 3.18 representa a distribuição das notas nas questões discursivas no componente de Conhecimento Específico.
Essa distribuição tem uma moda principal no intervalo [ 0 ; 10 ].
Nota -se que a partir desse intervalo há uma outra concentração de notas no intervalo ( 50 ; 60 ], também sendo relevante terem sido registradas notas em todos os intervalos.
O coeficiente de assimetria é negativo e próximo de zero ( -0,19 ).
Este coeficiente também é negativo nas regiões Sudeste, Sul e Centro-Oeste ( -0,21, -0,26 e -0,28, respectivamente ).
Já nas regiões Norte e Nordeste a assimetria foi à direita, com coeficiente positivo ( 0,08 e 0,09 respectivamente ). 3.3.2.1 Análise da Questão Discursiva 3 do Componente de Conhecimento Específico Na questão 3, cujos resultados aferidos encontram -se descritos na Tabela 3.12, a média dos estudantes de todo o Brasil foi 29,3.
A menor média nessa questão foi obtida pelos alunos da região Norte ( 23,6 ), enquanto a maior média foi obtida na região Centro-Oeste ( 30,4 ).
Quanto à variabilidade das notas, o desvio padrão de todo o Brasil foi 35,6.
O maior desvio padrão foi obtido na região Centro-Oeste ( 36,1 ), enquanto o menor foi obtido na região Norte ( 31,3 ).
A nota máxima, 100,0 pontos, foi alcançada por alunos de todas as regiões do país.
A mediana do Brasil como um todo foi 20,0, e foi a mesma nas regiões Sudeste, Sul e Centro-Oeste, enquanto nas demais regiões a mediana foi 0,0.
Ou seja, pelo menos metade dos estudantes das regiões Norte e Nordeste tirou nota zero nesta questão.
A mesma nota mínima ( 0,0 ) foi obtida em todas as regiões do Brasil, sem exceção.
O Gráfico 3.19 mostra a distribuição das notas na questão discursiva 3, do componente de Conhecimento Específico.
Essa distribuição é unimodal, com moda nas questões em branco.
Considerando -se em separado a distribuição das notas dos alunos que resolveram a questão 3, observamos outras concentrações nos intervalos ( 10 ; 20 ] e ( 90 ; 100 ].
O coeficiente de assimetria é positivo ( 0,89 ) para o Brasil como um todo e para as Grandes Regiões ( entre 0,81 na região Centro-Oeste e 1,22 na região Norte ), indicando uma assimetria à direita. 3.3.2.2 Comentários sobre as respostas à Questão Discursiva 3 A questão 3 apresentou enunciado claro e com nível de sofisticação adequado ao que se espera de graduandos.
O nível de profundidade exigido, inicialmente considerado intermediário, foi avaliado como avançado após as correções.
Apesar de bem construída e condizente com o conhecimento esperado para a Área de Gestão Comercial, a questão se apresentou com nível de dificuldade elevado para os alunos.
O enunciado pedia que o aluno realizasse uma ligação da teoria ( ambiente de marketing ) com a prática e propusesse uma inovação de produto coerente com a análise ambiental previamente desenvolvida.
Foi verificado que os alunos tiveram muita dificuldade nessa questão, principalmente na identificação do conteúdo teórico a ser tratado e na relação desse conteúdo à prática.
A questão que exigia conhecimento específico a respeito de uma das disciplinas do curso ( no caso, marketing ), e serviu para avaliar o que é fundamental para a gestão comercial.
As notas baixas indicam problemas na formação dos alunos, apesar da estratégia de pontuação ser bem subdividida.
As correções demonstraram que a estrutura de pensamento exigida para construir a resposta não parece ter sido desenvolvida pelos respondentes.
Alguns alunos confundiram características do ambiente de marketing com as características da organização orientada para marketing, identificando como características do ambiente sociocultural brasileiro : • “ engajamento de uma boa equipe de marketing ” ; • “ saber administrar os pontos fortes e fracos da empresa ” ; • “ analisar pontos fortes e pontos fracos da empresa - trabalhar e melhorar aquilo que é bom ”.
Outros respondentes foram evasivos na identificação das variáveis ambientais apontando, por exemplo : • “ A tecnologia cresce cada vez mais.
Está muito mais fácil se comunicar, manter relações pessoais e divulgar informações ” ; • “ A população brasileira da cada vez mais oportunidade ao comercio inovador ” ; • “ O ambiente sociocultural fica cada vez mais avançado ” ; Algumas respostas continham propostas de inovação muito evasivas : • “ ter um bom planejamento social, propaganda e comunicação ” ; • “ ter um produto de boa qualidade por um preço acessível ” ; • “ pensar na criação das alternativas para que o produto seja negociado ” ; • “ fazer pesquisa de mercado ”. 3.3.2.3 Análise da Questão Discursiva 4 do Componente de Conhecimento Específico A Tabela 3.13 contém as informações relativas à questão 4 do conjunto de questões do componente de Conhecimento Específico.
O desempenho dos estudantes de todo o Brasil nesta questão foi superior ao desempenho na questão de número 3.
A média geral do Brasil foi 53,4, sendo a menor média registrada na região Nordeste ( 48,2 ) e a maior na região Sul ( 54,7 ).
A nota máxima ( 100,0 ) foi atingida em todas as regiões.
A mediana do Brasil com um todo foi ( 70,0 ), sendo a mesma nas regiões Sudeste, Sul e Centro-Oeste, enquanto que nas regiões Norte e Nordeste a mediana foi 60,0.
A nota mínima ( 0,0 ) foi obtida em todas as regiões do Brasil, sem exceção.
O Gráfico 3.20, representa a distribuição de notas na questão discursiva 4, no componente de Conhecimento Específico.
Como nas outras questões discursivas, o número de estudantes que deixou a questão em branco foi muito grande, quase um quarto.
A distribuição é bimodal, com modas nas questões em branco e no intervalo ( 70 ; 80 ].
A assimetria da distribuição foi negativa ( -0,43 ), indicando uma assimetria à esquerda.
Na distribuição dos resultados dos alunos por região, os coeficientes de assimetria também são sempre negativos variando entre -0,52 na região Sul e -0,15 na região Nordeste. 3.3.2.4 Comentários sobre as respostas à Questão Discursiva 4 Como o enunciado não trazia informações suficientes, ao solicitar que fossem listadas possíveis causas para a queda nas vendas e que fossem propostas ações gerenciais, qualquer causa seria possível bem como qualquer proposta de ação gerencial.
Dessa forma, pelo padrão de resposta, os fatores externos à empresa que fossem influentes sobre o tipo de negócio apresentado poderiam ser : de cunho demográfico ( tais como idade, raça, gênero, escolaridade etc ) ou de cunho cultural ( tais como valores, crenças, estilo de vida etc ) ou, ainda, de cunho comportamental ( tais como moda, modismos, hábitos de consumo, mídias prestigiadas etc ).
O importante era sempre relacioná -los ao dono de animais de estimação e aos usuários de aparelhos móveis de acesso à internet.
Em algumas situações, o respondente utilizou informações ou conhecimentos que não estavam presentes no enunciado, provavelmente por ter algum interesse no tema, ou ter conhecimento do caso da Americanas.com.
Esse caso esteve presente na mídia algum tempo antes da realização das provas e deve ter servido como inspiração para responder à questão.
Outros respondentes pareceram confundir comércio eletrônico com a prestação de serviços de telefonia móvel, provavelmente por conta da proibição de comercialização de novas linhas imposta à operadora TIM, na mesma época da realização das provas.
Assim, a justificativa deveria sinalizar de que forma a variável apontada é relevante para o negócio.
Por exemplo, caso uma das variáveis escolhida fosse"posse de animais de estimação nos usuários de aparelhos móveis", a justificativa deveria indicar que essa informação será útil para definir o potencial de mercado ou a demanda esperada.
Ou ainda, se fosse escolhido"gênero dos usuários de aparelhos móveis", a justificativa poderia dizer que esse perfil define o tipo de conteúdo disponibilizado.
Foram frequentemente identificados problemas de grafia de palavras e de pontuação.
Por vezes, a leitura das respostas foi muito difícil e a compreensão comprometida.
Também se constatou que muitos respondentes confundiram causas e consequências, além de demonstrarem dificuldade de escrever de forma dissertativa. 3.3.2.5 Análise da Questão Discursiva 5 do Componente de Conhecimento Específico A Tabela 3.14 contém as informações relativas à questão 5 do conjunto do componente de Conhecimento Específico.
O desempenho dos estudantes nessa questão foi superior ao da questão 3 e inferior ao da questão 4.
A nota média dos estudantes de todo o Brasil foi 45,5.
A maior média foi registrada na região Sul ( 47,8 ), enquanto a menor média foi registrada na região Nordeste ( 38,4 ).
Quanto à variabilidade das notas, o desvio padrão dos alunos do Brasil, como um todo, foi 34,4.
Enquanto o maior desvio foi encontrado na região Nordeste ( 35,4 ), o menor foi encontrado na região Norte ( 32,4 ).
A maior nota máxima ( 100,0 ) foi alcançada na região Nordeste, enquanto a menor nota máxima ( 85,0 ) foi registrada nas regiões Norte, Sudeste e Centro-Oeste.
A mediana do país como um todo foi de 70,0, mesmo valor das regiões Sudeste, Sul e Centro-Oeste, enquanto as demais regiões tiveram mediana 40,0.
A nota mínima foi zero para todas as regiões do Brasil.
O Gráfico 3.21 representa a distribuição das notas da questão discursiva 5 do componente de Conhecimento Específico.
Destaca -se, como nas demais questões discursivas, o grande número de estudantes que deixaram esta questão em branco.
No entanto, pode -se identificar que a principal moda da distribuição está no intervalo ( 60 ; 70 ].
O coeficiente de assimetria para todos os alunos é -0,37 e para quase todas as regiões se mantém negativo ( entre -0,50 na região Sul e -0,21 na região Norte ).
A exceção é a região Nordeste cujo coeficiente é 0,01 com uma distribuição praticamente simétrica. 3.3.2.6 Comentários sobre as respostas à Questão Discursiva 5 A questão solicitava ao aluno que ele propusesse uma solução baseada em dois dos quatro princípios de negociação, mas, como isso é muito difícil, nenhuma resposta trouxe exatamente uma proposta desse tipo.
O uso de casos ilustrativos é desejável, mas no caso da questão 5 não havia uma relação clara entre o conhecimento que deveria ser utilizado e o contexto proposto.
Os alunos não conseguiram escolher dois critérios, como solicitado, e relacioná -los a uma solução para o problema enfrentado.
Os que responderam à questão conseguiram : correlacionar um critério a uma solução ou ainda conseguiram descrever dois critérios e ligá -los separadamente a duas soluções.
Apesar desse problema, notou -se que a maioria dos respondentes demonstrou que havia assimilado o conteúdo a respeito de negociação tratado na questão e sugeriram soluções satisfatórias para a resolução dos problemas.
Destaca -se que o padrão de resposta e os critérios de pontuação da questão 5 levaram em consideração os princípios, as justificativas e a proposta de solução para o caso. 3.3.3 Considerações Finais O processo de correção das questões discursivas do ENADE 2012 para os cursos superiores de Tecnólogos em Gestão Comercial revelou algumas dificuldades apresentadas pelos alunos.
Consideramos que a questão 3 era a mais clara e objetiva.
Essa questão avaliou bem o conhecimento dos alunos e foram percebidos poucos problemas de compreensão de seu enunciado.
A questão 4 exigia apenas bom senso para ser respondida, não cobrava conhecimento de um ponto específico ou de um conjunto de pontos disciplinares a serem discutidos.
Dessa forma, era possível responder à questão, sem precisar utilizar conhecimento específico do curso.
A questão 5 solicitava uma tarefa que era praticamente impossível devido ao tempo e tipo de avaliação.
O resultado foi muito ruim, já que não houve resposta que atendesse ao que foi solicitado.
As dificuldades apresentadas pelos alunos com o uso da língua portuguesa em sua forma escrita foram frequentes.
Além disso, foi percebida dificuldade de redigir um texto dissertativo que expressasse o pensamento do aluno acerca do que era solicitado.
Essas dificuldades indicam problemas graves na formação, mesmo que reconhecidamente oriundos da fase pré-universitária.

% matched lemmas: 
\end{textsample}
